\section{Conclusion}
\label{sec:conclusion}

Merge-Split MCMC~\citep{janson2022} narrows the stationary-distribution
ambiguity of standard \recom{}~\citep{deford2021} by sampling an independent
reverse spanning tree and computing the ratio
$\mathrm{cuts\_fwd} / \mathrm{cuts\_rev}$ as the Metropolis-Hastings
acceptance probability.
The forward cut count appears in the numerator because it dilutes the
forward proposal probability; the resulting ratio is the implemented
correction for the balanced-cut-count part of
$Q(\pi' \!\to\! \pi) / Q(\pi \!\to\! \pi')$.

The implementation as \texttt{MergeSplitChain} achieves $O(m \log m)$
per-step cost and raises the audit surface from ``always accept feasible
standard-ReCom proposals'' to a declared two-tree stochastic correction with
reverse zero-cut rejection.
The empirical acceptance rate of $\approx\!60\%$ on NC and WI is intermediate
between standard \recom{} ($\approx\!80\%$) and \forestrecom{} ($\approx\!40\%$),
reflecting the two-tree heuristic's intermediate strictness.

\mergesplit{} is a good production chain for reproducible redistricting
candidate search when the report can state its approximation boundary clearly.
\forestrecom{} is the escalation when the forest-measure claim is central;
SMC~\citep{dellaLibera2026smc} is the escalation when particle weights and
resampling diagnostics are the right calibration story.
The implementation is available via \texttt{bisect state --search merge-split}.
